\documentclass[12pt]{article}
\usepackage[utf8]{inputenc}
\usepackage[T1]{fontenc}
\usepackage{hyperref}
\title{Maridan Language Resources for K-12 Curriculum Development}
\author{Compiled by Claude Sonnet 5.5 (claude-sonnet-5-5) for Andr\'as Kornai}
\date{}
\begin{document}
\maketitle

\section{Introduction}
Maridan (ISO 639-3 zmd; Glottocode mari1422 in the survey data) is a Western Daly language of the Northern Territory, Australia. It is also recorded as Marridan, Meradan, Maredan, Murindan and Marrisyefin. Its country lay inland, along the marshlands north of the middle Moyle River. Endangered Languages, citing Moseley (2007), reports that in 1981 perhaps twenty speakers were reported and that only a few speakers are left. No literacy was recorded, and the searches found no school programme. Do not confuse it with the related Maridjabin (Marri Tjevin) or Marithiel, which are separate entries. Glottolog lists Maridan with the Marithielic subgroup.

\section{Primary Educational Resources}
No primers, textbooks or school programmes for Maridan turned up in the searches.

\section{Linguistic and Reference Materials}
\begin{itemize}
\item Wikipedia, Marridan: \url{https://en.wikipedia.org/wiki/Marridan}
\item Wikipedia, Marri Tjevin (related language, for disambiguation): \url{https://en.wikipedia.org/wiki/Marri_Tjevin}
\item Glottolog entry returned for Maridan (its code differs from the survey's, so verify): \url{https://glottolog.org/resource/languoid/id/mari1420}
\item Endangered Languages Project entry: \url{https://endangeredlanguages.com/lang/1907}
\item Handbook of the Cosmos (hgcosmos.org) Marridan page: \url{https://www.hgcosmos.org/assets/culture/AZ/M/Marridan.pdf}
\end{itemize}

\section{Religious and Community Resources}
The searches found no scripture translations or community organisations specific to Maridan.

\section{Digital and Audio Resources}
No audio or digital teaching resources were found in the searches. The Endangered Languages Project page is the place to look for any deposited recordings, but I did not confirm any.

\section{Curriculum Development Context}
Maridan is moribund or dormant, with at most a few remaining speakers and little documentation visible online. Any revitalization work would start with archival searching and consultation with the descendant community. The resources found here are overview and database entries only.

\end{document}